\documentclass{article}

% --- Codificación y idioma ---
\usepackage[utf8]{inputenc} % Codificación UTF-8
\usepackage[T1]{fontenc}
\usepackage[spanish]{babel}
\usepackage{lmodern}
\usepackage{microtype}

% --- Márgenes ---
\usepackage[margin=2.5cm]{geometry}

% --- Gráficos y tablas ---
\usepackage{graphicx}
\usepackage{booktabs}
\usepackage{caption}
\usepackage{subcaption}
\usepackage{siunitx}

% --- Matemáticas ---
\usepackage{amsmath,amssymb,amsthm}


\begin{document}

\title{Solución: Plano fase y Estabilidad}
\author{}
\date{}
\maketitle

\section{Problema}
Considere el siguiente sistema lineal numérico que modela un sólido con temperatura interna $X(t)$ en contacto con una ``capa'' o camisa $Y(t)$, estando el ambiente fijo en $A = 70$ (se trabaja con temperaturas relativas $x = X - A$, $y = Y - A$):

\begin{equation}
    \begin{cases}
        \dfrac{dx}{dt} = -4(x - y), \\
        \dfrac{dy}{dt} = -y
    \end{cases}
\end{equation}

\begin{enumerate}
    \item Calcule los puntos críticos y clasifíquelos (tipo y estabilidad).
    \item Verifíquelo construyendo el plano de fase y explique brevemente el significado físico de las trayectorias.
\end{enumerate}

\section{Solución}

\subsection{Puntos críticos}

Para encontrar los puntos críticos, igualamos a cero ambas ecuaciones:

\begin{align}
    \frac{dx}{dt} &= -4(x - y) = 0 \label{eq:dx} \\
    \frac{dy}{dt} &= -y = 0 \label{eq:dy}
\end{align}

De la ecuación (\ref{eq:dy}): $-y = 0 \implies y = 0$

Sustituyendo en (\ref{eq:dx}): $-4(x - 0) = 0 \implies -4x = 0 \implies x = 0$

Por lo tanto, el único punto crítico es:
\[
(x, y) = (0, 0)
\]

\subsection{Clasificación del punto crítico}

Escribimos el sistema en forma matricial:
\[
\begin{pmatrix}
\dot{x} \\ \dot{y}
\end{pmatrix}
=
\begin{pmatrix}
-4 & 4 \\
0 & -1
\end{pmatrix}
\begin{pmatrix}
x \\ y
\end{pmatrix}
\]

Calculamos los autovalores de la matriz:
\[
A = \begin{pmatrix}
-4 & 4 \\
0 & -1
\end{pmatrix}
\]

El polinomio característico es:
\[
\det(A - \lambda I) = \det\begin{pmatrix}
-4-\lambda & 4 \\
0 & -1-\lambda
\end{pmatrix} = (-4-\lambda)(-1-\lambda) - (4)(0)
\]
\[
= (\lambda + 4)(\lambda + 1) = 0
\]

Los autovalores son:
\[
\lambda_1 = -4, \quad \lambda_2 = -1
\]

\begin{align}
    (\text{tr } A)^{2} &> 4 |A| \\
    25 &> 16
\end{align}


\textbf{Clasificación:}
\begin{itemize}
    \item Ambos autovalores son reales y negativos
    \item $\lambda_1 = -4 < 0$, $\lambda_2 = -1 < 0$ , raíces de la ecuación característica.
    \item Tipo de equilibrio: \textbf{Nodo estable} , \[(tr A)^2 > 4 |A| \]
    \item Estabilidad: \textbf{Asintóticamente estable}
\end{itemize}

\subsection{Significado físico}

El punto crítico $(0, 0)$ corresponde a:
\[
x = X - 70 = 0 \implies X = 70
\]
\[
y = Y - 70 = 0 \implies Y = 70
\]

Esto significa que el sistema alcanza el equilibrio térmico cuando tanto el sólido como la capa tienen la misma temperatura del ambiente ($70^\circ$). Las trayectorias representan cómo las temperaturas relativas $x$ y $y$ decaen exponencialmente hacia cero, indicando que el sistema tiende al equilibrio térmico con el ambiente.

\end{document}